% ===================== 第 6 章 备考计划 =====================
\section{备考计划（2026 年 12 月考试）}

以 2026 年 12 月考试（预计 12 月中旬）倒推，共约十周。核心节奏：\textbf{前期背词 + 分块突破，中期真题 + 错题复盘，后期模考 + 保手感}。

\subsection{十周总表}

\begin{center}
\small
\begin{longtable}{p{1.6cm}p{3.1cm}p{5.3cm}p{4.6cm}}
\toprule
\textbf{周次} & \textbf{时间} & \textbf{词汇任务} & \textbf{真题/专项任务} \\
\midrule
\endfirsthead
\toprule
\textbf{周次} & \textbf{时间} & \textbf{词汇任务} & \textbf{真题/专项任务} \\
\midrule
\endhead
W1 & 9.29--10.5 & 第 2 章动词 1--75 & 1 套听力 + 1 篇仔细阅读 \\
\midrule
W2 & 10.6--10.12 & 第 2 章动词 76--150 & 2 篇仔细阅读 + 1 篇长篇阅读 \\
\midrule
W3 & 10.13--10.19 & 第 2 章形容词 & 写作 1 篇（现象解释型）+ 翻译 2 篇 \\
\midrule
W4 & 10.20--10.26 & 第 2 章名词 & 写作 1 篇（谚语型）+ 翻译 2 篇 \\
\midrule
W5 & 10.27--11.2 & 第 2 章副词 + 易混词 & 选词填空专项 + 1 套阅读 \\
\midrule
W6 & 11.3--11.9 & 第 3 章话题词汇（教育/科技/环境） & 完整套题 1（不含作文） \\
\midrule
W7 & 11.10--11.16 & 第 3 章话题词汇（经济/文化/健康） & 完整套题 2（含作文 1 篇） \\
\midrule
W8 & 11.17--11.23 & 第 3 章话题词汇（网络/成长/职业/交通） & 完整套题 3 + 错题复盘 \\
\midrule
W9 & 11.24--11.30 & 全量滚动复习一遍 & 完整套题 4（限时 130 分钟） \\
\midrule
W10 & 12.1--12.7 & 只背不熟的“顽固词” & 完整套题 5 + 三篇翻译重做 \\
\midrule
W11 & 12.8--考前 & 高频短语与写译句型 & 模考 1 套 + 放松调整，不碰新题 \\
\bottomrule
\end{longtable}
\end{center}

\subsection{每日模板（可选其一）}

\begin{multicols}{2}
\textbf{工作日（约 1.5 小时）}
\begin{itemize}
  \item 晨间 20 分钟：单词一组（约 40 个），遮住释义自测
  \item 午间 30 分钟：1 篇仔细阅读或 1 段听力
  \item 晚间 40 分钟：隔日轮换“写作/翻译”，写一句算一句
\end{itemize}

\textbf{周末（约 3 小时）}
\begin{itemize}
  \item 限时做半套或整套真题
  \item 逐题复盘：为什么错、考点在哪、下次怎么避
  \item 把新词补进自己的生词本，标注来源年份
\end{itemize}
\end{multicols}

\begin{tipbox}[title=复盘比刷题更重要]
每套真题做 1 遍、复盘 2 遍。复盘时回答三个问题：\\
(1) 这道题的考点是什么（词义 / 搭配 / 同义替换 / 逻辑）？\quad
(2) 我为什么错（词不认识 / 定位错 / 推断过度）？\quad
(3) 下次遇到同类题，我的第一步动作是什么？
\end{tipbox}

\subsection{考前一周 Checklist}

\begin{itemize}
  \item 打印或导出近三年真题各 1 套，按考试时间（15:00--17:25）完整模拟一次。
  \item 背熟写作“四类框架”各自的首段、过渡、结尾模板句各 3 句。
  \item 默写翻译高频专有名词表（节气、朝代、四大名著、非遗）。
  \item 检查耳机、调频频率、准考证与文具；提前熟悉考点交通。
  \item 考前一晚不熬夜；听力状态靠睡出来的。
\end{itemize}
